\chapter{Desenvolvimento}
\label{cap:desenvolvimento}

Este capítulo descreve a implementação realizada a partir do projeto do \autoref{cap:proj}. Inicialmente são apresentados o ambiente e a organização do código, incluindo a reprodução da implementação original de \citeonline{Souza2023}, que é pré-requisito de todas as comparações. Na sequência, é descrito o fluxo de execução, do carregamento da cena às métricas. Em seguida, são apresentados os outros três eixos de otimização --- paralelismo, automação de hiperparâmetros e segmentadores ---, com destaque para a estratégia de fusão concorrente proposta pelo orientador. O capítulo termina com a organização dos experimentos e os desvios em relação ao projeto.

\section{Ambiente e Organização do Código}
\label{sec:ambiente}

\subsection{Ambiente de Execução}

A implementação foi desenvolvida em Python 3.12, com PyTorch 2.13 (CUDA 13.0) para a rede neural \cite{Paszke2019PyTorch}, scikit-learn para o PCA e as métricas, OpenCV para os algoritmos de segmentação e Optuna para a Otimização Bayesiana \cite{Akiba2019Optuna}. Os experimentos foram executados em um computador pessoal com GPU NVIDIA RTX 5070 Laptop, processador com 32 núcleos lógicos e 15,8 GB de memória RAM, sob WSL2. O ambiente virtual foi criado com o gerenciador \texttt{uv}, e as versões exatas de todas as bibliotecas ficam fixadas em um arquivo de requisitos versionado junto ao código e são conferidas ao iniciar o notebook.

\subsection{Organização em Pacote}

O trabalho de referência foi entregue como um notebook Jupyter, com estado global compartilhado entre células e configurações espalhadas pelo código. Para que cada resultado deste trabalho fosse atribuível a uma configuração conhecida, a implementação foi reorganizada em um pacote Python, com um módulo por responsabilidade e cada definição em um único lugar. O \autoref{qua:pacote} resume a organização. O notebook foi mantido, mas apenas importa o pacote e guarda a narrativa e as células de execução; os experimentos também integram o pacote, de modo que a implementação e as medições compartilham exatamente o mesmo código.

\begin{quadro}[H]
    \centering
    \caption{Organização do pacote da implementação}
    \begin{tabular}{|p{0.22\textwidth}|p{0.70\textwidth}|}
        \hline
        \rowcolor[HTML]{C0C0C0}
        \textbf{Módulo} & \textbf{Responsabilidade} \\
        \hline
        Configuração e ambiente & Parâmetros de uma execução em um único objeto imutável; seleção de dispositivo, sementes e modo determinístico \\
        \hline
        Dados & Carregamento das cenas, tabela de pixels, PCA, composição RGB, os quatro segmentadores e a divisão em treino, validação e teste \\
        \hline
        Modelo, treino e avaliação & CNN 1D com a fusão, laço de treino com parada antecipada e cálculo de OA, AA e Kappa \\
        \hline
        Pipeline & A única peça que conhece a ordem das etapas, chamada pelo notebook, pela busca de hiperparâmetros e pelos experimentos \\
        \hline
        Paralelismo & Execução concorrente dos ramos do pré-processamento, partição da cena e fusão concorrente \\
        \hline
        Hiperparâmetros & Espaço de busca, as três estratégias de busca, seleção dos vencedores e campanha final \\
        \hline
        MLOps & Registro imutável de cada execução, campanhas e relatórios \\
        \hline
        Visualização & Diagnósticos, explicabilidade e as figuras deste trabalho \\
        \hline
        Experimentos & Um módulo por experimento, com retomada a partir dos resultados já gravados \\
        \hline
    \end{tabular}
    \captionsetup{justification=raggedright}
    \caption*{Fonte: Elaborado pelo autor}
    \label{qua:pacote}
\end{quadro}

\subsection{Reprodução da Implementação Original}
\label{subsec:reproducao}

O notebook original foi escrito em TensorFlow/Keras, e a reimplementação foi feita em PyTorch. Antes de qualquer otimização, foi necessário garantir que a reimplementação fosse fiel ao original e que suas execuções se repetissem; sem isso, nenhuma diferença medida depois poderia ser atribuída à otimização. A reimplementação preserva a arquitetura do notebook --- 37.577 parâmetros em University of Pavia, idênticos aos do resumo do modelo Keras ---, o PCA com branqueamento, a divisão dos dados com as mesmas sementes do notebook (semente 42, com as mesmas 3.208 amostras de teste em University of Pavia) e os hiperparâmetros de treino: taxa de aprendizado 0,001, \textit{dropout} 0,5, lote 32 e teto de 60 épocas. A regra de parada antecipada foi trocada, como descrito na \autoref{subsec:desTreino}. Dois padrões do Keras precisaram ser replicados explicitamente, porque os do PyTorch são outros: a inicialização Xavier uniforme com viés nulo nas convoluções e nas camadas densas e o $\epsilon = 10^{-7}$ do otimizador Adam. O notebook original não fixa a semente da inicialização e do treino da rede, e a execução registrada nele obteve OA 0,9819, AA 0,9735 e Kappa 0,9760 em University of Pavia. Como os dois frameworks não partem dos mesmos pesos, a igualdade numérica com o Keras não é esperada: a equivalência é de arquitetura, de dados e de hiperparâmetros. A reimplementação, por sua vez, é determinística: com a configuração do notebook, execuções repetidas produzem o mesmo histórico de perda por época, os mesmos pesos finais e as mesmas métricas (OA 0,976309, AA 0,957918 e Kappa 0,968545).

A reprodutibilidade das execuções seguintes depende de três mecanismos. O primeiro é a configuração explícita: todos os parâmetros de uma execução ficam em um único objeto, que é gravado junto ao resultado e permite reconstruí-la. O segundo é o modo determinístico, que fixa as sementes do NumPy e do PyTorch, desativa a escolha automática de algoritmos do cuDNN e a precisão reduzida TF32, e restringe o OpenCV a uma única thread. Esta última restrição foi necessária porque o LSC paralelo do OpenCV devolvia rótulos diferentes a cada chamada, como discutido na \autoref{sec:resOrganizacao}. O terceiro é um contrato de dados verificado antes de cada treino: os conjuntos de treino, validação e teste precisam ser disjuntos, os espectros não podem conter valores inválidos, as classes precisam estar numeradas de forma contígua e as duas entradas da rede precisam ter uma linha por amostra.

Duas práticas de engenharia completam o eixo de organização. Cada \textit{trial} da busca de hiperparâmetros mantinha sua própria cópia dos dados preparados, embora dele só interessassem as métricas: 2,5 MB por \textit{trial} em Indian Pines, ou cerca de 450 MB em uma campanha de 180 \textit{trials}. Adotou-se, então, o ciclo de instanciar, treinar, extrair as métricas e liberar os recursos de cada \textit{trial}, de modo que apenas o melhor resultado corrente permanece em memória; com ele, a memória do processo cresceu 16 MB ao longo de uma campanha desse tamanho (\autoref{sec:resOrganizacao}). Além disso, toda execução longa tornou-se retomável: os resultados são gravados a cada configuração concluída de um experimento, e uma execução interrompida continua de onde parou, relendo o que já foi medido.

\section{Fluxo de Execução}
\label{sec:fluxo}

Uma execução completa percorre as etapas a seguir, sempre na mesma ordem:

\begin{enumerate}
    \item \textbf{Carregamento da cena}: o cubo hiperespectral $(H, W, B)$ e o mapa de verdade terrestre $(H, W)$ são lidos e achatados em uma tabela de $H \cdot W$ pixels;
    \item \textbf{Ramo espectral}: os pixels sem rótulo (\textit{background}) são removidos e o PCA com branqueamento (\textit{whitening}) reduz as $B$ bandas a $K_b$ componentes (30 em Indian Pines e 15 em University of Pavia e Salinas Valley, como no original);
    \item \textbf{Ramo espacial}: um PCA de 3 componentes sobre a cena inteira, incluindo o \textit{background}, gera uma composição RGB de 8 bits, que o segmentador divide em superpixels; o identificador da região de cada pixel rotulado é a segunda entrada da rede;
    \item \textbf{Divisão dos dados}: as amostras são divididas em treino, validação e teste e, opcionalmente, o treino é subamostrado;
    \item \textbf{Treino}: a CNN 1D é construída e treinada com parada antecipada;
    \item \textbf{Avaliação}: o modelo treinado classifica o conjunto de teste, e OA, AA, Kappa e a matriz de confusão são calculados.
\end{enumerate}

Os dois ramos são independentes até o ponto de fusão. É essa independência que as estratégias de paralelismo da \autoref{sec:desParalelismo} exploram.

\subsection{Segmentadores}
\label{subsec:desSegmentadores}

Os quatro segmentadores --- SEEDS, SLIC (no modo SLICO), LSC e Watershed --- são implementados sobre o OpenCV e devolvem saídas no mesmo formato: o mapa de rótulos $(H, W)$ e a máscara de contornos. Por isso, trocar o segmentador não altera nenhuma outra etapa. O Watershed do OpenCV não fornece a máscara de contornos, que é então derivada do mapa de rótulos, e recebe marcadores dispostos em uma grade regular; as linhas que ele marca entre bacias foram mantidas como uma região própria.

Os parâmetros de cada algoritmo são definidos por cena (\autoref{qua:paramSeg}). Indian Pines usa regiões menores, porque tem apenas 145 $\times$ 145 pixels e classes agrícolas em faixas estreitas (\autoref{fig:indian_pines}). O número de superpixels pedido ao SEEDS é um alvo, não uma garantia: em Indian Pines, 350 superpixels pedidos resultam em 81 regiões; em University of Pavia, 150 resultam em 45.

\begin{quadro}[H]
    \centering
    \caption{Parâmetros padrão dos segmentadores por cena}
    \begin{tabular}{|p{0.16\textwidth}|p{0.38\textwidth}|p{0.36\textwidth}|}
        \hline
        \rowcolor[HTML]{C0C0C0}
        \textbf{Algoritmo} & \textbf{Indian Pines} & \textbf{University of Pavia e}\newline\textbf{Salinas Valley} \\
        \hline
        SEEDS & 350 superpixels, 200 iterações, \textit{prior} 2, 10 níveis, 5 \textit{bins} de histograma & 150 superpixels, 150 iterações, \textit{prior} 1, 8 níveis, 15 \textit{bins} \\
        \hline
        SLIC & região de 8 px, 10 iterações & região de 12 px, 10 iterações \\
        \hline
        LSC & região de 8 px, razão 0,075, 10 iterações & região de 12 px, razão 0,075, 10 iterações \\
        \hline
        Watershed & 350 marcadores & 500 marcadores \\
        \hline
    \end{tabular}
    \captionsetup{justification=raggedright}
    \caption*{Fonte: Elaborado pelo autor. Os valores do SEEDS são os do notebook original de \citeonline{Souza2023}; o texto daquele trabalho relata 150 superpixels em Indian Pines e, em University of Pavia e Salinas Valley, \textit{prior} 2 e 3 níveis}
    \label{qua:paramSeg}
\end{quadro}

\subsection{Modelo e Fusão}
\label{subsec:desModelo}

A CNN 1D segue o notebook original e está detalhada no \autoref{qua:camadas} para Indian Pines. As quatro convoluções não usam preenchimento (\textit{padding}), e cada uma encurta a sequência em $k - 1$ posições, consumindo 14 posições no total; por isso a rede exige pelo menos 15 componentes espectrais. Em University of Pavia e Salinas Valley, com 15 componentes, a última convolução produz uma única posição, e o achatamento resulta em 64 características.

\begin{quadro}[H]
    \centering
    \caption{Camadas da CNN 1D com fusão, em Indian Pines ($K_b = 30$)}
    \begin{tabular}{|p{0.36\textwidth}|p{0.30\textwidth}|p{0.18\textwidth}|}
        \hline
        \rowcolor[HTML]{C0C0C0}
        \textbf{Camada} & \textbf{Saída}\newline\textbf{(canais, comprimento)} & \textbf{Parâmetros} \\
        \hline
        Entrada: espectro após PCA & (1, 30) & --- \\
        \hline
        Conv1D 16 filtros, $k = 9$, ReLU & (16, 22) & 160 \\
        \hline
        Conv1D 32 filtros, $k = 3$, ReLU & (32, 20) & 1.568 \\
        \hline
        Conv1D 64 filtros, $k = 3$, ReLU & (64, 18) & 6.208 \\
        \hline
        Conv1D 64 filtros, $k = 3$, ReLU & (64, 16) & 12.352 \\
        \hline
        Achatamento & 1.024 & --- \\
        \hline
        Fusão: concatenação do identificador do superpixel & 1.025 & --- \\
        \hline
        Densa 128, ReLU, \textit{dropout} & 128 & 131.328 \\
        \hline
        Densa 64, ReLU, \textit{dropout} & 64 & 8.256 \\
        \hline
        Densa de saída (\textit{logits}) & 16 & 1.040 \\
        \hline
        \textbf{Total} & & \textbf{160.912} \\
        \hline
    \end{tabular}
    \captionsetup{justification=raggedright}
    \caption*{Fonte: Elaborado pelo autor}
    \label{qua:camadas}
\end{quadro}

A fusão é tardia: as convoluções nunca veem o identificador do superpixel, que só é concatenado ao vetor de características depois do achatamento, como um número a mais. A cabeça densa é a única parte da rede que enxerga as duas informações juntas. Na implementação, a concatenação foi separada da cabeça densa em um passo próprio, o que permite à fusão concorrente da \autoref{subsec:fusaoConcorrente} proteger apenas esse passo com uma trava. Como alternativa à concatenação do identificador numérico, também foi implementada uma tabela de vetores aprendidos por região (\textit{embedding}), usada apenas para entender como o identificador é aproveitado pela rede; a configuração herdada, com o identificador numérico, é a referência de todas as comparações.

\subsection{Treino, Parada Antecipada e Divisão dos Dados}
\label{subsec:desTreino}

O treino usa a perda de entropia cruzada e o otimizador Adam. A parada antecipada segue uma política explícita: o treino é interrompido após 10 épocas sem que a acurácia de validação melhore em pelo menos 0,001, e os pesos da melhor época são restaurados. Nos experimentos de segmentadores e de hiperparâmetros (05, 06, 08 e 10; \autoref{qua:experimentos}), há ainda uma meta: o treino é interrompido quando a acurácia de validação atinge 0,95, com os pesos da melhor época igualmente restaurados, e o teto é de 200 épocas (no experimento 08, orçamentos de 10, 40 e 200 épocas por \textit{trial}). Nos experimentos de partição (11 e 13) e na fusão concorrente (14), vale apenas a paciência, com teto de 25 e de 60 épocas, respectivamente; o teto de 60 épocas é o do notebook original. O original, porém, parava pela perda de validação, com paciência de 5 épocas; a troca para a acurácia de validação foi feita para que a mesma métrica governe a parada e a seleção de hiperparâmetros. Em todos os casos, a escolha de modelos e de hiperparâmetros usa exclusivamente o conjunto de validação, e as métricas de teste são apenas relatadas.

A divisão dos dados reproduz a do original: um sorteio estratificado por classe separa 85\% das amostras para treino e divide os 15\% restantes igualmente entre validação e teste. Opcionalmente, apenas uma fração do treino é usada, com amostragem estratificada; validação e teste permanecem intactos, para que as métricas continuem comparáveis entre frações. Os experimentos de segmentadores e de hiperparâmetros (05, 06, 08 e 10) usam 10\% do treino, tanto na busca quanto no retreino do vencedor; os demais usam o treino completo.

\section{Paralelismo}
\label{sec:desParalelismo}

Três estratégias de paralelismo foram implementadas, em níveis diferentes do fluxo. A primeira executa concorrentemente os dois ramos do pré-processamento; a segunda divide a cena em partes processadas ao mesmo tempo; a terceira sobrepõe a segmentação à extração de características da rede. O paralelismo muda quando cada passo é executado, e não o que ele calcula, e os experimentos de partição e de fusão concorrente conferem isso: na partição da cena, a OA das execuções em processos foi igual à da execução sequencial das mesmas partições em todos os valores de $N$ medidos; na fusão concorrente, as 723 execuções do experimento 14 previram as mesmas classes que a inferência em um único lote (\autoref{subsec:resFusao}). A exceção é o treino independente das partições em threads: as partições compartilham o gerador de números aleatórios do PyTorch, consumido em ordem imprevisível na inicialização dos pesos e nas máscaras de \textit{dropout}, e execuções idênticas variaram até 3,51 pontos de OA (\autoref{sec:resOrganizacao}). Com processos, cada partição tem o seu gerador, e o resultado se repete.

\subsection{Entre os Ramos do Pré-processamento}

A preparação dos dados pode executar os ramos espectral e espacial de três formas: em sequência, que é a linha de base; em duas threads, que compartilham os \textit{arrays} de entrada sem cópia; ou com o ramo espacial em um processo filho, que exige serializar a tabela de pixels e os resultados. A versão com threads funciona apesar do GIL porque a maior parte do processamento dos dois ramos ocorre em código nativo que libera o GIL: a álgebra linear do PCA e a iteração do segmentador no OpenCV. Pelo segundo corolário da \autoref{sec:concorrencia}, o ganho máximo é dado pelos tempos dos dois ramos, e o ramo espacial, que contém a segmentação, é o mais longo, com 78\% a 93\% do tempo sequencial do pré-processamento (\autoref{qua:ramos}).

\subsection{Por Partição da Cena}

Para paralelizar a etapa mais custosa do fluxo, o treino, a cena é dividida em $N$ partições com número semelhante de pixels rotulados, de modo que o teto teórico do \textit{speedup} se aproxime de $N$. A segmentação é executada uma única vez sobre a cena inteira, antes da partição: segmentar cada parte separadamente produziria superpixels diferentes nas fronteiras, e os identificadores deixariam de ser comparáveis entre partições. As partições são faixas horizontais contíguas, e dois modos de uso foram implementados: o \textit{replicado}, em que um único modelo é treinado com a cena inteira e apenas a inferência é paralelizada, e o \textit{independente}, em que cada partição treina seu próprio modelo com as amostras dela. Nos experimentos de partição, o treino usa lote 64, em vez dos 32 da receita herdada, e o \textit{speedup} é medido contra a execução sequencial das mesmas partições, no mesmo dispositivo. Como o lote é o mesmo nos dois lados da comparação, o \textit{speedup} compara executores sob a mesma carga, mas as acurácias desses experimentos não são comparáveis às da linha de base.

As partições podem ser executadas em sequência, em threads ou em processos. Os processos são criados pelo método \textit{spawn}, que inicia um novo interpretador, sem herdar o estado do processo pai, porque o PyTorch interrompe com erro a cópia de processo (\textit{fork}) feita depois que o cálculo de gradientes foi executado em threads \cite{PyTorchMultiprocessing}. O custo do \textit{spawn} é alto, porque cada processo reimporta o PyTorch: com 16 partições, a partida chega a 78\% do tempo de uma execução com o pool criado dentro da medição (\autoref{qua:poolQuente}). Por isso, foi implementado um pool reutilizável: os processos são criados e aquecidos uma única vez, fora da medição, e atendem às execuções seguintes. Cada processo é configurado para usar uma única thread do PyTorch, pois, do contrário, $N$ processos com 16 threads cada disputariam os mesmos núcleos. Antes de criar um pool, a memória necessária é estimada a partir da memória ocupada por processo, cerca de 481 MB, e um pool que não caberia na RAM é recusado em vez de levar o sistema a usar a área de troca.

\subsection{Fusão Concorrente}
\label{subsec:fusaoConcorrente}

O projeto previa que a segmentação e a CNN 1D fossem executadas ao mesmo tempo, com sincronização apenas no ponto de fusão (\autoref{fig:paralelizacao}). A estratégia de paralelismo entre os ramos, porém, executou concorrentemente à segmentação apenas o PCA espectral; a CNN só era executada depois. O orientador deste trabalho detalhou, então, um esquema que cumpre a proposta original, representado na \autoref{fig:fusaoConcorrente}:

\begin{itemize}
    \item uma thread T0 executa o ramo espacial --- PCA de 3 componentes, composição RGB e SEEDS --- e publica o mapa de segmentação (MS) em uma estrutura compartilhada;
    \item um pool de threads T1 a TN recebe as posições do mapa: cada tarefa passa o espectro dos pixels daquelas posições apenas pelas convoluções da rede, lê o identificador do segmento de cada posição no mapa e faz a concatenação, e a cabeça densa produz o vetor de classes;
    \item a segmentação e a concatenação são protegidas por um mutex.
\end{itemize}

\begin{figure}[H]
    \centering
    \captionsetup{justification=centering}
    \caption{Fusão concorrente: T0 segmenta enquanto o pool T1 a TN extrai as características e realiza a fusão}
    \begin{tikzpicture}[
        font=\footnotesize,
        caixa/.style={draw, rounded corners=2pt, minimum height=0.85cm, align=center, inner sep=4pt},
        thread/.style={draw, dashed, rounded corners=4pt, inner sep=6pt},
        seta/.style={-{Stealth[length=2mm]}, thick}
    ]
        \node[caixa] (hsi) {HSI\\$(H, W, B)$};
        \node[caixa, right=0.7cm of hsi] (pca3) {PCA(3)\\RGB};
        \node[caixa, right=0.7cm of pca3] (seeds) {SEEDS};
        \node[caixa, right=0.7cm of seeds, fill=gray!15] (ms) {Mapa de\\segmentação (MS)};
        \node[thread, fit=(pca3)(seeds), label={[font=\footnotesize\bfseries]above:T0}] (t0) {};

        \node[caixa, below=1.5cm of hsi] (pix) {Espectro do\\pixel (PCA)};
        \node[caixa, right=0.7cm of pix] (conv) {CNN 1D\\(só as convoluções)};
        \node[caixa, right=0.7cm of conv] (vet) {Vetor de\\características};
        \node[circle, draw, right=1.4cm of vet, inner sep=1pt] (cat) {$+$};
        \node[caixa, right=0.7cm of cat] (dens) {Camadas\\densas};
        \node[caixa, right=0.6cm of dens] (cls) {Vetor de\\classes};
        \node[thread, fit=(conv)(vet)(cat), label={[font=\footnotesize\bfseries]below:T1 a TN (uma posição do MS por tarefa)}] (tn) {};

        \draw[seta] (hsi) -- (pca3);
        \draw[seta] (pca3) -- (seeds);
        \draw[seta] (seeds) -- (ms);
        \draw[seta] (pix) -- (conv);
        \draw[seta] (conv) -- (vet);
        \draw[seta] (vet) -- (cat);
        \draw[seta] (ms.east) -| node[pos=0.75, right, align=left]{id do segmento\\na posição $(i, j)$} (cat.north);
        \draw[seta] (cat) -- (dens);
        \draw[seta] (dens) -- (cls);
    \end{tikzpicture}
    \captionsetup{justification=raggedright}
    \caption*{Fonte: Elaborada pelo autor, a partir do esquema proposto pelo orientador}
    \label{fig:fusaoConcorrente}
\end{figure}

As convoluções não dependem do mapa de segmentação, e é isso que permite às threads de trabalho avançar enquanto T0 ainda segmenta: só a concatenação precisa esperar. Duas decisões de implementação foram necessárias para transformar o esquema em um programa correto.

A primeira diz respeito à ordem. Um mutex garante que T0 e as threads de trabalho não acessem o mapa ao mesmo tempo, mas não garante que T0 o acesse primeiro: se uma thread de trabalho adquirir a trava antes de T0, lerá um mapa ainda vazio. Por isso, a trava do esquema foi implementada como a de uma variável de condição (\autoref{sec:concorrencia}), cujo predicado é um indicador de que o mapa já foi publicado. Uma thread que chega antes da publicação aguarda, liberando a trava durante a espera, e T0, ao publicar o mapa, notifica todas as threads bloqueadas. Para tornar visível qualquer leitura indevida, o mapa começa preenchido com um valor que nenhum segmentador produz.

A segunda diz respeito à extensão da seção crítica. Foram implementadas duas versões, comparáveis porque todo o resto é idêntico:

\begin{itemize}
    \item \textbf{como no esquema}: T0 detém a trava durante toda a segmentação, e cada tarefa a detém durante a concatenação;
    \item \textbf{mínima}: T0 adquire a trava apenas para publicar o mapa pronto, e cada tarefa, apenas para ler os identificadores. Depois de publicado, o mapa passa a ser somente lido, e leituras concorrentes de um dado imutável não produzem condição de corrida; manter a concatenação dentro da trava serializa as threads sem proteger nada.
\end{itemize}

A granularidade das tarefas também é um parâmetro, medido com blocos de 1, 64 e 1.024 posições. Com um bloco de uma posição, cada tarefa recebe um único pixel, exatamente como no esquema; com blocos maiores, cada tarefa recebe um conjunto contíguo de posições, o que dilui o custo fixo de despachar uma tarefa e de chamar a rede. Esse custo pesa: em Indian Pines, chamar a rede pixel a pixel, mesmo sem threads, já torna a inferência 3,12 vezes mais lenta que a de todos os pixels em um único lote, e com uma posição por tarefa o esquema fica de 7 a 35 vezes mais lento (\autoref{subsec:resFusao}).

Na implementação do esquema, cada tarefa faz as convoluções e a fusão em sequência; é a forma chamada de \textit{como no esquema} no \autoref{qua:fusao}, na \autoref{fig:fusaoSpeedup} e na \autoref{fig:fusaoLinhaTempo}. Isso tem uma consequência estrutural: uma thread que termina as convoluções antes de T0 publicar o mapa fica bloqueada à espera dele e, enquanto isso, não atende nenhuma outra tarefa da fila. Enquanto T0 segmenta, portanto, no máximo $N$ tarefas avançam, e a sobreposição pretendida fica limitada a $N$ vezes o tamanho do bloco. Para medir o teto do esquema, foi implementada também uma \textit{variante desacoplada}, em que convolução e fusão são tarefas separadas e todas as convoluções entram na fila antes de qualquer fusão. Como a fila do pool é atendida em ordem de chegada, quando uma fusão é retirada da fila, a convolução correspondente já foi iniciada, e nenhuma thread fica ociosa enquanto ainda há convolução a fazer. As duas extensões da seção crítica foram medidas nas duas formas.

Três cuidados completam a implementação. O modo de cálculo de gradientes do PyTorch é definido por thread; por isso, a inferência sem gradientes precisa ser ativada dentro de cada tarefa, e não apenas na thread principal. Pelo mesmo motivo, o número de threads internas do PyTorch, que nas medições de paralelismo é restrito a uma, precisa ser fixado em cada thread do pool: a configuração feita na thread principal não se propaga, e as threads do pool usariam as 32 threads padrão do OpenMP, o que fazia a convolução de um bloco de 64 pixels levar cerca de 100 ms em vez de 0,2 ms. Por fim, as threads do pool são criadas sob demanda e inicializam estado interno do PyTorch na primeira operação; o pool é, então, aquecido antes das medições com uma barreira que obriga a criação de todas as threads, separando o custo de partida do custo do paralelismo, como no pool de processos.

O esquema se aplica à propagação direta (\textit{forward pass}) de um modelo já treinado, isto é, à inferência. No treino, a perda de cada lote depende do identificador do segmento desde o primeiro lote, e o gradiente atravessa as convoluções; além disso, a segmentação é executada uma única vez. Não há, portanto, trabalho de treino que possa ser sobreposto à segmentação.

\section{Automação de Hiperparâmetros}
\label{sec:desHPO}

O espaço de busca combina os hiperparâmetros de treino, o número de componentes do PCA, o segmentador e três receitas de parâmetros de cada segmentador (\autoref{qua:espaco}). Os valores de cada hiperparâmetro incluem o da receita herdada, de modo que a configuração original é um dos candidatos possíveis. Em Indian Pines o espaço tem $4 \times 3 \times 3^4 \times 3 = 2.916$ combinações. Com o tempo médio de treino de um \textit{trial} do experimento 08 com teto de 200 épocas (3,2 s, com 10\% do treino), percorrê-las levaria cerca de 2,6 h de treino por semente, 146 vezes o custo de uma busca de 20 \textit{trials}, e a comparação das estratégias repete cada busca com 3 sementes. Por isso, a busca em grade prevista no projeto foi substituída pela busca aleatória, que, assim como a grade, não usa os resultados anteriores para escolher o próximo candidato e é a referência usual para avaliar estratégias adaptativas \cite{Bergstra2012Random}.

\begin{quadro}[!htbp]
    \centering
    \caption{Espaço de busca de hiperparâmetros}
    \begin{tabular}{|p{0.28\textwidth}|>{\RaggedRight\arraybackslash}p{0.40\textwidth}|>{\RaggedRight\arraybackslash}p{0.20\textwidth}|}
        \hline
        \rowcolor[HTML]{C0C0C0}
        \textbf{Hiperparâmetro} & \textbf{Valores} & \textbf{Herdado} \\
        \hline
        Taxa de aprendizado & 0,0003; 0,001; 0,003 & 0,001 \\
        \hline
        \textit{Dropout} & 0,20; 0,35; 0,50 & 0,50 \\
        \hline
        \textit{Weight decay} & 0; $10^{-5}$; $10^{-4}$ & 0 \\
        \hline
        Tamanho do lote & 32; 64; 128 & 32 \\
        \hline
        Componentes do PCA & 0,5$\times$, 1$\times$ e 1,5$\times$ o valor da cena, com mínimo de 15 (Indian Pines: 15, 30, 45; University of Pavia e Salinas Valley: 15, 22) & 30 (Indian Pines), 15 (University of Pavia e Salinas Valley) \\
        \hline
        Segmentador & SEEDS, SLIC, LSC, Watershed & SEEDS \\
        \hline
        Receita do segmentador & a padrão e duas variações em torno dela: no SEEDS, 0,7$\times$ e 1,35$\times$ o número de superpixels, com 120 e 250 iterações e \textit{prior} 1 e 3; no SLIC, o tamanho da região $\pm$3 px; no LSC, o tamanho da região $\pm$3 px, com razão 0,04 e 0,12; no Watershed, 0,7$\times$ e 1,35$\times$ o número de marcadores & a padrão \\
        \hline
    \end{tabular}
    \captionsetup{justification=raggedright}
    \caption*{Fonte: Elaborado pelo autor}
    \label{qua:espaco}
\end{quadro}

Todo candidato é representado da mesma forma: uma lista de índices, um por hiperparâmetro. Essa representação é comum às três estratégias implementadas, o que garante que elas difiram apenas na forma de escolher o próximo candidato:

\begin{itemize}
    \item \textbf{busca aleatória}: sorteia os índices de forma uniforme, sem repetir candidatos;
    \item \textbf{TPE}: usa o amostrador da biblioteca Optuna, tratando segmentador e receita como hiperparâmetros categóricos e os demais como inteiros ordenados \cite{Bergstra2011TPE, Akiba2019Optuna};
    \item \textbf{algoritmo genético}: população de 6 candidatos no experimento de comparação das estratégias, seleção por torneio entre 3, cruzamento uniforme, mutação de cada gene com probabilidade de 0,25 e elitismo dos 2 melhores, que passam à geração seguinte sem novo treino \cite{Holland1975}.
\end{itemize}

As três estratégias começam com a mesma cobertura: os quatro primeiros \textit{trials} são obrigados a usar segmentadores distintos, para que nenhum fique de fora da busca por acaso. Cada \textit{trial} é avaliado pela maior acurácia de validação de seu histórico. Na busca de baixa fidelidade, a paciência da parada antecipada passa a ser um terço do orçamento de épocas quando esse valor é menor que 10, para que ainda possa ser acionada: com 10 épocas por \textit{trial}, ela cai para 3; com 40 e com 200 épocas, permanece 10. Ao final da busca, apenas o vencedor é retreinado, com teto de 200 épocas e 5 sementes (experimentos 06 e 10), e a decisão --- o que foi buscado, o que venceu e o plano final --- é gravada em arquivos com \textit{hash}, de modo que seja possível justificar posteriormente a promoção de cada receita.

\section{Segmentadores sob Otimização Conjunta}
\label{sec:desSegmentadores}

O plano de verificação previa testar os segmentadores mantendo a CNN fixa. Esse desenho foi implementado e descartado (\autoref{cap:resultados}): a configuração da CNN que serve a um segmentador não serve a outro, e a comparação com a CNN fixa mede quanto uma receita arbitrária combina com cada algoritmo, e não a qualidade do algoritmo. No uso real, a CNN e o segmentador são ajustados juntos, e a comparação foi refeita dessa forma: para cada segmentador, uma busca TPE de 16 \textit{trials}, restrita a ele, cujo vencedor, escolhido pela acurácia de validação, é retreinado com 5 sementes. O SEEDS passa pelo mesmo tratamento e é a referência.

\section{Organização dos Experimentos}
\label{sec:desExperimentos}

A campanha planejada no \autoref{tab:descricaoExperimentos} foi desdobrada em experimentos menores, cada um com um objetivo específico, conforme o \autoref{qua:experimentos}. A numeração segue a ordem em que foram criados; os números ausentes correspondem a experimentos absorvidos por outros, descartados quando o escopo foi restrito à CNN 1D, ou não abordados nesta monografia. Cada experimento é executável a partir de um lançador de linha de comando, grava seus resultados em arquivos tabulares versionados e confere, quando aplicável, a repetibilidade das próprias medições.

\begin{quadro}[H]
    \centering
    \caption{Experimentos realizados e correspondência com o plano de verificação}
    \begin{tabular}{|p{0.07\textwidth}|p{0.46\textwidth}|p{0.19\textwidth}|p{0.16\textwidth}|}
        \hline
        \rowcolor[HTML]{C0C0C0}
        \textbf{Exp.} & \textbf{Objetivo} & \textbf{Eixo} & \textbf{Plano} \\
        \hline
        05 & Comparação dos segmentadores com a CNN fixa & Segmentadores & Exp. 2 (contraexemplo) \\
        \hline
        06 & Comparação dos segmentadores, cada um com hiperparâmetros otimizados & Segmentadores & Exp. 2 e 4 \\
        \hline
        08 & Comparação entre busca aleatória, TPE e algoritmo genético; efeito do orçamento de épocas; importância de cada hiperparâmetro & Hiperparâmetros & Exp. 3 e 4 \\
        \hline
        10 & Receita otimizada em comparação com a herdada & Hiperparâmetros & Exp. 1, 3 e 4 \\
        \hline
        11 & Escalabilidade da partição da cena, nos modos replicado e independente & Paralelismo & Exp. 5 \\
        \hline
        13 & Pool frio em comparação com pool quente & Paralelismo & Exp. 5 \\
        \hline
        14 & Fusão concorrente & Paralelismo & Exp. 5 \\
        \hline
    \end{tabular}
    \captionsetup{justification=raggedright}
    \caption*{Fonte: Elaborado pelo autor}
    \label{qua:experimentos}
\end{quadro}

\section{Desvios em Relação ao Projeto}
\label{sec:desvios}

A implementação manteve os requisitos do projeto, mas algumas escolhas mudaram ao longo do desenvolvimento. O \autoref{qua:desvios} reúne as que dizem respeito à implementação, e o \autoref{qua:desviosExperimentos}, as que dizem respeito à campanha experimental, cada uma com o seu motivo.

\begin{quadro}[H]
    \centering
    \caption{Desvios em relação ao projeto: implementação}
    \begin{tabular}{|>{\RaggedRight\arraybackslash}p{0.27\textwidth}|>{\RaggedRight\arraybackslash}p{0.27\textwidth}|>{\RaggedRight\arraybackslash}p{0.36\textwidth}|}
        \hline
        \rowcolor[HTML]{C0C0C0}
        \textbf{Previsto} & \textbf{Realizado} & \textbf{Motivo} \\
        \hline
        TensorFlow/Keras ou PyTorch & PyTorch & Arquitetura, PCA, divisão dos dados e hiperparâmetros de treino do original preservados em PyTorch, com execução determinística; a regra de parada antecipada foi trocada (última linha deste quadro) \\
        \hline
        CNN com 16, 32, 64 e 128 filtros, \textit{max-pooling} e \textit{softmax} & 16, 32, 64 e 64 filtros, sem \textit{pooling}, saída em \textit{logits} & A arquitetura do notebook original, que é a referência; o \textit{softmax} é aplicado pela função de perda \\
        \hline
        Rótulo do superpixel codificado em \textit{one-hot} & Identificador numérico concatenado, como no original & A configuração herdada é a referência; o \textit{embedding} foi implementado para análise \\
        \hline
        Critério do RF03: classificação com fusão superior à puramente espectral & Fusão verificada pela concatenação e pela dimensão da entrada da cabeça densa (1.025 em Indian Pines); a vantagem sobre a rede puramente espectral é a reportada por \citeonline{Souza2023} & A fusão faz parte da configuração de referência e não é variável de estudo; nenhuma configuração sem fusão é treinada \\
        \hline
        SEEDS, LSC e Watershed & SEEDS, SLIC, LSC e Watershed & O SLIC é uma referência frequente nas avaliações de algoritmos de superpixels \cite{Stutz2018} e uma comparação natural \\
        \hline
        Busca em grade e Otimização Bayesiana com Ray Tune & Busca aleatória, TPE (Optuna) e algoritmo genético & As 2.916 combinações custariam 146 vezes uma busca de 20 \textit{trials} (\autoref{sec:desHPO}); o genético acrescenta uma estratégia populacional \\
        \hline
        Épocas, filtros, tamanho dos \textit{kernels} e otimizador entre os hiperparâmetros da CNN & Taxa de aprendizado, \textit{dropout}, \textit{weight decay} e lote; arquitetura e otimizador Adam fixos & A arquitetura e o otimizador são os do original, que é a referência; o orçamento de épocas foi estudado à parte, no experimento 08 \\
        \hline
        Google Colab ou HPC & GPU local (RTX 5070 Laptop) & Suficiente para as três cenas, com controle total do ambiente; os 15,8 GB de RAM impediram o pool de 32 processos (\autoref{sec:resParalelismo}) \\
        \hline
        Parada do notebook original: perda de validação, paciência 5 & Acurácia de validação, paciência 10, melhores pesos restaurados; meta de 0,95 nos experimentos de segmentadores e de hiperparâmetros & A mesma métrica governa a parada e a seleção \\
        \hline
    \end{tabular}
    \captionsetup{justification=raggedright}
    \caption*{Fonte: Elaborado pelo autor}
    \label{qua:desvios}
\end{quadro}

\begin{quadro}[H]
    \centering
    \caption{Desvios em relação ao projeto: campanha experimental}
    \begin{tabular}{|>{\RaggedRight\arraybackslash}p{0.27\textwidth}|>{\RaggedRight\arraybackslash}p{0.27\textwidth}|>{\RaggedRight\arraybackslash}p{0.36\textwidth}|}
        \hline
        \rowcolor[HTML]{C0C0C0}
        \textbf{Previsto} & \textbf{Realizado} & \textbf{Motivo} \\
        \hline
        Segmentadores com CNN fixa & Segmentadores sob otimização conjunta & O desenho com CNN fixa se mostrou inválido (\autoref{cap:resultados}) \\
        \hline
        85\% das amostras para treino & 10\% desse treino nos experimentos de segmentadores e de hiperparâmetros & Cada execução fica cerca de 5,7 vezes mais barata em Indian Pines (\autoref{qua:fracaoTreino}), o que viabilizou os 540 \textit{trials} do experimento 08; validação e teste permanecem inteiros \\
        \hline
        Todos os experimentos nas três cenas & CNN fixa (05) em Indian Pines e University of Pavia; segmentadores (06) e fusão concorrente (14) nas três; busca (08), ganho da busca de hiperparâmetros (10) e partição (11 e 13) em Indian Pines & Custo de tempo; Indian Pines é a cena de menor acurácia na linha de base (OA 0,9051; \autoref{qua:linhaBase}) \\
        \hline
        Testes de significância estatística & Diferença afirmável quando excede duas vezes o desvio combinado das duas condições entre sementes & Com 3 a 5 sementes por condição, um critério simples e explícito; por usar o desvio das execuções, e não o erro padrão da média, é mais exigente que o mesmo limiar aplicado ao erro padrão; o pareamento por semente não foi explorado \\
        \hline
        Segmentação em processo e CNN na GPU, em paralelo & Threads e processos em CPU, em três níveis & Em Indian Pines, o pré-processamento é 0,9\% de uma execução e a inferência, 0,4\% (\autoref{sec:resParalelismo}); o tempo concentra-se no treino, e a comparação entre executores exige o mesmo dispositivo \\
        \hline
        Paralelização sobre a melhor configuração identificada & Fusão concorrente sobre a configuração herdada; partição com lote 64 e 25 épocas & Interessa o tempo entre executores, com a mesma configuração dos dois lados \\
        \hline
    \end{tabular}
    \captionsetup{justification=raggedright}
    \caption*{Fonte: Elaborado pelo autor}
    \label{qua:desviosExperimentos}
\end{quadro}
